\section{Introduction}

\subsection{Motivation}
Systems programming in C remains foundational to modern computing
infrastructure\autocite{dist_systems_languages, C_embedded_programming}.
Operating systems, hardware drivers, and performance-critical software
all depend on the C language's direct access to
memory\autocite{memory_errors_stridharan}.
Yet the same flexibility that makes C efficient also makes it fragile:
small inconsistencies in layout can cause catastrophic runtime
errors \autocite{evaluation_memory_errors, C_errors_study}.
Layout-related faults—incorrect padding, mis-aligned fields, or unintended
aliasing—manifest as memory corruption and undefined behaviour.
These faults are notoriously difficult to detect, especially when they arise
from hand-crafted structures that must precisely conform to hardware
or protocol formats.

In response, verified programming languages such as
\textbf{Cogent}\autocite{dargent}
have have demonstrated that a functional, type-safe systems code
can be compiled to efficient C with formal correctness guarantees.
However, Cogent and similar languages abstract away and exist in their own
ecosystems. The \textbf{Dargent} framework withing Cogent partly addresses this
issue by allowing the specification of \textit{data layouts} alongside
high-level data types, generating C structs and accessors that obey those
specifications.

Despite its promise, Dargent exposes an inherent challenge: layouts must be
low-level enough to control bits and bytes, yet compositional enough
to reason and verify. Its design is also closely tied to Cogent's compiler,
which makes reuse outside that context difficult. These limitations motiviate
the development of an independent, declarative language for describing and
generating correct low-level layouts—a language that can serve both as a
research vehicle for reasoning about memory representations, and as a
practical tool for C programmers. 

\subsection{Research Aim and Question}
This thesis explores how a small, declarativce language can express complex
memory layouts while preserving reasoning about correctness. The project
defines a prototype langauge, \textbf{Cedar}, originally proposed by
Anderson\autocite{cedar_george}, which treats data layouts as a first-class
concern. Cedar allows programmers to describe how values of structured types
occupy memory, decoupling layout from logical structure.

The central \textbf{research question} guiding this work is:
\begin{quote}
    \textbf{How can the Cedar language be designed and evaluated to support
    practical, verifiable specification of data layouts in C?}
\end{quote}

Answering this question requires both language-design research and
implementation engineering. The goal is not only to reproduce Dargent-style
behaviour in an independent compiler, but also to generalise and clarify
underlying principles: compositional layout descriptions, well-typed accessors,
and generation of C code.

\subsection {Approach and Methodology}
The research follows a \textit{design-and-evaluation} methodology typical
of programming-language research\autocite{incremental_development}.
The work proceeds in three interlocking stages:
\begin{enumerate}
    \item \textbf{Language Design and Formalisation — } analyising the previous
    Cedar proposal and revising its intermediate representations
    (L $\rightarrow$ LR $\rightarrow$ CR) to better capture nexsted records, arrays,
    and bitfields.
    \item \textbf{Implementation and Code Generation — } implementing the Cedar
    compiler to emit cohesive, self-contained C output for representative
    examples, enabling direct comparison with Dargent's generated code.
    \item \textbf{Evaluation via Case Studies — } running the Cedar compiler on
    examples to assess expressivity, correctness, and consistency of generated
    C output.
\end{enumerate}

Throughout these stages, the design is informed by principles from declarative
programming, as it is implemented in a functional style using Haskell.
Rather aiming for a complete verification effort, this thesis focuses
\textit{enabling correctness}: defining interfaces, type rules, and translation
structures that would later support formal reasoning. 
    
\subsection{Contributions}
The main contributions of this thesis are:
\begin{itemize}
    \item \textbf{An extended Language Design for Cedar —} supporting nexsted
    record layouts, bitfields, and more expressive alignment semantics, while
    retaining declarative composition.
    \item \textbf{A Refined Compilation Pipeline —} a modular transofromation
    from surface syntax to layout-resolved forms
    ( L $\rightarrow$ LR $\rightarrow$ CR ) culminating in generated C code that
    preserves structure and alignment.
    \item \textbf{A Prototype Implementation —} a working Cedar compiler
    implemented in Haskell, integrating parsing, type checking layout matching,
    and code generation.
    \item \textbf{Empirical Evaluation via Case Studies —}  reproduction
    of Dargent examples to demonstrate Cedar's usability and internal
    consistency.
    \item \textbf{Discussion of Design Trade-offs and Future Verification
    Pathways —} identifying the boundaries between declarative specification,
    compiler automation, and machine-checkable proof of correctness.
\end{itemize}

Collectively, these outcomes advance Cedar from a concept to a more
complete research language, and contribute new insights into the design of
trustworthy data-layout languages for C.

\subsection{Thesis Structure}
The remainder of this thesis is organised as follows:
\begin{itemize}
    \item \textbf{Chapter 2} reviews background material on C memory layouts,
    verified compilation, the Cogent and Dargent frameworks, and the origins
    of Cedar.
    \item \textbf{Chapter 3} describes the research methodology and language design,
    detailing Cedar's intermediate representations and translation
    pipeline.
    \item \textbf{Chapter 4} presents the implementation and evaluation,
    including case case studies. 
    \item \textbf{Chapter 5} discusses design implications, limitations,
    and directions for future verification and usability improvements.
    \item \textbf{Chapter 6} concludes with a summary of findings and prospects
    for integrating Cedar into broader verified-systems toolchains.
    
\end{itemize}


